\subsection{\textit{Sprint} 4}

A última \textit{sprint} reuniu ajustes do aplicativo e funcionalidades de
busca, moderação e denúncias. A equipe trabalhou nas telas administrativas,
na consulta de usuários e estabelecimentos, no tratamento de denúncias e na
gestão de categorias. Também havia pendências de configuração para o uso do
aplicativo fora do ambiente de desenvolvimento, como domínio da \ac{api},
notificações e links de convite. Concluímos as atividades previstas para a
etapa final.

No \textit{backend}, implementei a exclusão administrativa de
\textit{tags} e categorias. A operação verificava se a categoria ainda
continha outras \textit{tags} e se o item estava associado a eventos ou
usuários. Para uma \textit{tag} em uso, a remoção exigia confirmação explícita.
Também era registrado quem havia realizado a ação, mantendo o histórico da
alteração para auditoria.

Trabalhei ainda nos links de convite para eventos. No \textit{backend},
preparei a página de destino e os arquivos que associavam o domínio ao
aplicativo. No \textit{frontend}, completei o suporte à abertura desses links
no iOS e no Android, por meio de \textit{universal links} e
\textit{app links} \cite{apple-universal-links,android-app-links}.

Essas tarefas complementavam o compartilhamento de convites que já existia.
O link precisava funcionar ao ser enviado por outros aplicativos e conduzir
o usuário ao evento. Com o Hangy instalado, a abertura deveria ocorrer no
próprio aplicativo. Sem uma sessão ativa, o fluxo preservava o convite para
retomá-lo depois do \textit{login}. A página de destino atendia a quem
acessasse o link pelo navegador.

A equipe também incluiu uma tarefa de segurança com Trivy
\cite{trivy-docs}, voltada a verificar vulnerabilidades e impedir a
implantação quando fossem identificadas falhas de gravidade alta ou crítica.

A semana de provas e as demais demandas de fim de semestre limitaram o
tempo disponível. Precisei conciliar as implementações com as
responsabilidades de gestão e com as avaliações. Nesse período, a atenção
às dependências continuava necessária, pois uma configuração de domínio ou
do aplicativo podia impedir o uso de uma funcionalidade já implementada.

Ao verificar os convites, precisei conferir o link compartilhado, a página
servida pelo \textit{backend} e a configuração nativa do aplicativo. Também
considerei se o usuário tinha o aplicativo instalado e uma sessão ativa
para avaliar o funcionamento do fluxo completo.
